\section{Related work}

\paragraph{Advantage design and sample polarity.}
\citet{zhu2025nsr} report better high-$k$ pass@$k$ retention under negative sample reinforcement (NSR) than positive sample reinforcement (PSR). Subsequent work examines negative-sample exploration \citep{tang2025polarity} and asymmetric or difficulty-dependent advantages \citep{ingle2026negative,yu2026symmetry,decugis2026fade}. The original NSR--PSR comparison motivates our matched-pass@1 test on two model families. On Qwen, trace-matched zero-mean sign noise added to PSR reproduces NSR's pooled early accuracy lag, with seed heterogeneity. This establishes a perturbation that can reproduce the lag without separately identifying scalar noise magnitude and covariance structure.

\paragraph{The objective view of RLVR.}
Binary-reward updates can be expressed as ascent on monotone transforms of per-problem success \citep{davis2025objective}. This objective view also appears in surrogate-reward interpretations of advantage shaping \citep{thrampoulidis2025shaping} and analyses of weights that minimize time to target accuracy \citep{heckel2026weighting}. Pass@$k$ policy gradients are reweighted pass@1 gradients \citep{yu2025diagnostic}. Our extension is an exact per-problem second moment for correctness-only rules under independent sampling, allowing reward-dependent stopping and unequal conditional mean squared score-function norms $F_0$ and $F_1$. Shared parameters allow changes in prompt weights to rotate aggregate updates \citep{barakat2026interference}, so the mean identity does not imply a common frontier.

\paragraph{Diversity, pass@$k$, and evaluation.}
DARS already plots pass@1--pass@$K$ trajectories and reports higher pass@128 at fixed pass@1, without intervals or a registered matching protocol \citep{yang2025depth}. We add pre-registered decision criteria and problem-level bootstrap intervals that recompute the interpolated reference within shared pass@1 support. Base--RLVR pass@$k$ crossovers \citep{yue2025limit} and confidence bands for them \citep{yang2026kernel} make evaluation uncertainty relevant to these comparisons. CoT-pass@$k$ instead evaluates intermediate reasoning \citep{wen2025cotpassk}; BBG gates updates by posterior marginal pass@$k$ utility \citep{yuan2026overtraining}. At matched pass@1, pass@$k$ summarizes the cross-problem distribution of success probabilities. Our model predicts aggregate pass@16 using base success counts, observed checkpoint pass@1, and a transfer coefficient frozen before computing test predictions across rules, model families, and evaluation budgets. Its registered warm-start failure limits this description, while ``RLVR is a Kernel'' challenges deterministic prediction from base success more generally \citep{yang2026kernel}.

\paragraph{Replay, off-policy guidance, and rollout allocation.}
TAR and EliteReplay replay diverse base-model solutions during RL \citep{li2026replay,moussa2026elitereplay}. LSPO scaffolds all-failure groups, LUFFY incorporates off-policy demonstrations, and ExGRPO prioritizes stored rollouts \citep{ding2026lspo,yan2025guidance,zhan2025exgrpo}. HORA maximizes posterior hit utility under a fixed budget with per-problem normalization \citep{wang2026hora}; Knapsack RL reallocates exploration budgets \citep{knapsack2025}. DARS adds rollouts to difficult problems, increasing its budget relative to its baseline \citep{yang2025depth}, and AR3PO combines allocation with response reuse \citep{zhang2025reuse}. MaxRL interpolates expected-reward and maximum-likelihood objectives as sampling compute grows \citep{tajwar2026maxrl}. We evaluate replay, base-seeded replay and larger groups at matched pass@1, and use a longer reference run where an intervention leaves the 150-step reference range.

\paragraph{Exploration bonuses and entropy.}
Unlikeliness reward upweights rare correct trajectories \citep{he2025unlikely}, while entropy-based controls regulate exploration through token-level updates \citep{cui2025entropy}. These interventions depend on trajectory likelihood or token entropy, extending beyond correctness-only advantage rules. We measure entropy separately on fixed text and policy-generated responses, and compare training checkpoints with a late-checkpoint temperature sweep. Higher temperature does not exceed the training frontier at matched pass@1 in our setting.
